% Table V — Computational Complexity Analysis (P2 Task 1)
% N: number of agents; d: feature dim; k: avg. neighbors (k ≪ N under sparse R_c);
% T: temporal memory horizon (GRU step is O(Nd) per timestep).
%
% Note: pairwise distance for graph construction is O(N^2), shared by GAT/DSGF.
% Message-passing / attention complexity is what differs under radius constraints.

\begin{table}[!t]
\centering
\caption{Per-timestep computational complexity of guidance (or policy) modules.
$N$ denotes the number of UAVs, $d$ the hidden dimension, $k$ the average number of
neighbors within the communication radius ($k\ll N$ under sparse connectivity),
and $T$ the temporal memory horizon.}
\label{tab:complexity}
\begin{tabular}{lccc}
\toprule
\textbf{Method} & \textbf{Spatial Aggregation} & \textbf{Temporal} & \textbf{Complexity} \\
\midrule
MAPPO & --- & --- & $O(Nd)$ \\
GAT & Dense / radius graph & --- & $O(N^{2}d)$ \\
Transformer & Full attention & --- & $O(N^{2}d)$ \\
DSGF (Ours) & Quality-weighted sparse graph & GRU memory & $O(Nkd + Nd)$ \\
\bottomrule
\end{tabular}
\end{table}

% Compact CSV mirror for scripts / non-LaTeX readers:
% Method,Spatial,Temporal,Complexity
% MAPPO,-,-,O(Nd)
% GAT,dense graph,-,O(N^2 d)
% Transformer,attention,-,O(N^2 d)
% DSGF,sparse quality graph,GRU,O(Nkd+Nd)
